\documentclass[12pt]{article}
\usepackage[utf8]{inputenc}
\usepackage[russian]{babel}
\usepackage{amsmath}
\usepackage{amssymb}
\usepackage{geometry}
\geometry{a4paper, margin=2cm}

\begin{document}

\section*{Задача 3.1.6a}

\textbf{Условие:} Отсортировать последовательность \(a_1, \dots, a_n\) в порядке возрастания методом пузырька.\\
in: n, a1, ..., an \\
out: b1, ..., bn

\section*{Псевдокод}
Я использую сортировку пузырьком-обменом с лекций
\begin{verbatim}
for i:=2 to n-1
    for j:=n downto i
        if a[j-1].key > a[j].key then 
            a[j-1] :=: a[j]
\end{verbatim}

\section*{Решение}

\begin{verbatim}
Read 1      ; n
Load =1
Store 8     ; idx := 1

READ_LOOP: Load 8
    Sub 1
    JGTZ INIT_SORT
    Load 8
    Add =9
    Store 9
    Read *9
    Load 8
    Add =1
    Store 8
    Jump READ_LOOP
INIT_SORT: Load =2
    Store 2     ; i := 2
OUTER: Load 1
    Sub 2
    JGTZ INNER_INIT
    JZero INNER_INIT
    Jump PRINT_INIT
INNER_INIT: Load 1
    Store 3     ; j := n
INNER: Load 3
    Sub 2
    JGTZ CHECK
    JZero CHECK
    Jump NEXT_I
CHECK: Load 3
    Add =9
    Store 4     ; r4 := 9+j
    Load *4
    Store 7     ; r7 := a[j]
    Load 4
    Sub =1
    Store 5     ; r5 := 9+j-1
    Load *5
    Store 6     ; r6 := a[j-1]
    Load 6
    Sub 7
    JGTZ SWAP
    Jump NO_SWAP
SWAP: Load 6
    Store *4    ; a[j] := a[j-1]
    Load 7
    Store *5    ; a[j-1] := a[j]
NO_SWAP: Load 3
    Sub =1
    Store 3     ; j := j-1
    Jump INNER
NEXT_I: Load 2
    Add =1
    Store 2     ; i := i+1
    Jump OUTER
PRINT_INIT: Load =1
    Store 8
PRINT_LOOP: Load 8
    Sub 1
    JGTZ END
    Load 8
    Add =9
    Store 9
    Write *9
    Load 8
    Add =1
    Store 8
    Jump PRINT_LOOP
END: Halt 0
\end{verbatim}

\section*{Сложность}

Худший случай - массив, отсортированный по убыванию, тогда на каждой итерации внутреннего цикла выполняется swap.

\begin{enumerate}
    \item Инициализация - 3 шага
    \item Цикл READ\_LOOP выполняется n раз. Каждая итерация занимает 11 шагов, выход из цикла - 3 шага. Итого: 11n +3
    \item Инициализация сортировки - 2 шага
    \item Внешний цикл OUTER: $i = 2 \dots n$ - n-1 итераций. Внутренний цикл INNER: $j = n \dots i$ - n-i+1 итераций
    \begin{itemize}
        \item Для \(j > i\) каждая итерация со swap - 24 шага
        \item Для \(j = i\) - это последняя итерация внутреннего цикла - 25 шагов
        \item Выход из внутреннего цикла - 5 шагов, переход к следующему i - 4 шага
    \end{itemize}
    Тогда время внешнего цикла:
    \[
        S = \sum_{i=2}^{n} \bigl[24(n-i) + 25 + 5 + 4\bigr] = \sum_{i=2}^{n} \bigl[24(n-i) + 34\bigr] = 12n^2 - 2n - 10.
    \]
    \item При i=n+1 - выход из цикла 5 шагов
    \item PRINT\_INIT 2 шага
    \item PRINT\_LOOP там точно также как в READ\_LOOP, 11n + 3
    \item Halt 0 - 1 шаг
\end{enumerate}

Значит сложность в худшем случае:
\[
    T_{\max} = 3 + (11n + 3) + 2 + (12n^2 - 2n - 10) + 5 + 2 + (11n + 3) + 1 = 12n^2 + 20n + 9
\]

\end{document}